\subsection{Depth and acyclicity}

PROPOSITION 3.--- Let A be a ring, C a bounded-below complex of A-modules, and p an integer. Suppose that for every pair of integers \((m,n)\) with \(m \geqslant n \geqslant p\) , the depth of the A-module \(C_{m}\) relative to the annihilator of \(\mathrm{H}_{n}(\mathrm{C})\) is \textgreater m - n. We then have \(\mathrm{H}_{n}(\mathrm{C}) = 0\) for \(n \geqslant p\) .

Since C is bounded on the left, \(\mathrm{H}_{n}(\mathrm{C})\) is zero for n sufficiently large. If the conclusion were false, there would exist an integer \(q \geqslant p\) such that \(\mathrm{H}_{n}(\mathrm{C}) = 0\) for n \textgreater{} q and \(\mathrm{H}_{q}(\mathrm{C}) \neq 0\) . Let J denote the annihilator of \(\mathrm{H}_{q}(\mathrm{C})\) ; we then have \(\operatorname{prof}_{\mathrm{A}}(\mathrm{J}; \mathrm{H}_{q}(\mathrm{C})) = 0\) . Moreover, since \(\mathrm{Z}_{q}(\mathrm{C})\) is a submodule of \(C_{q}\) , and since by hypothesis \(\operatorname{prof}_{\mathrm{A}}(\mathrm{J}; \mathrm{C}_{q}) > q - q = 0\) , one has \(\operatorname{prof}_{\mathrm{A}}(\mathrm{J}; \mathrm{Z}_{q}(\mathrm{C})) > 0\) . We then deduce from the exact sequence

\[
0 \to \mathrm{B} _ {q} (\mathrm{C}) \to \mathrm{Z} _ {q} (\mathrm{C}) \to \mathrm{H} _ {q} (\mathrm{C}) \to 0
\]

the equality \(\mathrm{prof}_{\mathrm{A}}(\mathrm{J};\mathrm{B}_q(\mathrm{C})) = 1\) (n° 1, prop. 1). By the definition of \(q\) , \(\mathrm{B}_n(\mathrm{C})\) is equal to \(\mathrm{Z}_n(\mathrm{C})\) for every integer \(n > q\) . From the canonical exact sequences

\[
0 \to \mathrm{B} _ {n} (\mathrm{C}) \to \mathrm{C} _ {n} \to \mathrm{B} _ {n - 1} (\mathrm{C}) \to 0 \quad (n > q)
\]

and from the hypothesis \(\operatorname{prof}_{\mathrm{A}}(\mathrm{J};\mathrm{C}_{n}) > n - q\) , we obtain by induction the equality \(\operatorname{prof}_{\mathrm{A}}(\mathrm{J};\mathrm{B}_{n}(\mathrm{C})) = n - q + 1\) for every \(n \geqslant q\) (loc. cit.). But this is absurd, since \(\mathrm{B}_{n}(\mathrm{C})\) is zero for \(n\) sufficiently large.

COROLLARY 1.--- Let A be a ring, J an ideal of A, C a bounded-below complex of A-modules, and p an integer. Suppose that \(\mathrm{JH}_{m}(\mathrm{C}) = 0\) and \(\operatorname{prof}_{\mathrm{A}}(\mathrm{J}; \mathrm{C}_{m}) > m - p\) for \(m \geqslant p\) . One then has \(\mathrm{H}_{n}(\mathrm{C}) = 0\) for \(n \geqslant p\) .

Indeed, for \(n \geqslant p\) the annihilator \(J_n\) of \(H_n(C)\) contains \(J\) ; therefore we have \(\mathrm{prof}_{\mathrm{A}}(\mathrm{J}_{n};\mathrm{C}_{m}) \geqslant \mathrm{prof}_{\mathrm{A}}(\mathrm{J};\mathrm{C}_{m})\) ( \(n^{\circ}\) 1, cor. 1 of prop. 2), so that the hypothesis of the proposition is satisfied.

COROLLARY 2.-- Let A be a local ring, C a bounded-below complex of A-modules, p an integer. Suppose that for \(m \geqslant p\) , \(\mathrm{H}_{m}(\mathrm{C})\) is of finite length and \(\mathrm{C}_{m}\) of depth \(>m-p\) . One then has \(\mathrm{H}_{n}(\mathrm{C}) = 0\) for \(n \geqslant p\) .

The A-module \(\bigoplus_{m\geqslant p}\mathrm{H}_{m}(\mathrm{C})\) is of finite length. Let J denote its annihilator; by A, VIII, § 1, no. 3, corollary, the ring A/J is Artinian, hence J contains a power of the maximal ideal of A (A, VIII, § 10, no. 1, th. 1). Consequently one has \(\mathrm{prof}_A(\mathrm{J};\mathrm{C}_m)\geqslant \mathrm{prof}(\mathrm{C}_m) > m - p\) for \(m\geqslant p\) (no. 1, cor. 1 of prop. 2), so that cor. 1 can be applied.

